\section{Part 1 --- Profiling Execution Time}

\subsection{Overview}

We profiled 5 important functions in our voice agent pipeline using Python's \texttt{cProfile} module:

\begin{table}[h]
\centering
\begin{tabular}{|c|l|l|l|}
\hline
\textbf{\#} & \textbf{Function} & \textbf{Component} & \textbf{What it does} \\
\hline
1 & \texttt{VAD.\_process\_audio\_frame} & Voice Activity Detection & Runs Silero VAD inference on each 20ms audio chunk \\
2 & \texttt{ASR.\_transcribe\_audio} & Speech-to-Text & Sends audio to Groq Whisper API for transcription \\
3 & \texttt{LLM} TTFT via \texttt{litellm.completion} & Language Model & Calls Groq LLM API and streams tokens back \\
4 & \texttt{TTS} TTFB via \texttt{requests.post} & Text-to-Speech & Calls Inworld TTS API and streams audio back \\
5 & \texttt{Channel.send} / \texttt{\_wait\_and\_get} & Inter-thread messaging & Passes data (frames) between pipeline components \\
\hline
\end{tabular}
\caption{The 5 profiled functions.}
\end{table}

Functions 1--4 form the audio pipeline that determines \textbf{Time to First Audio (TTFA)} --- the latency from end-of-speech to first agent audio. Function 5 is the glue that connects them.

We first profile each audio pipeline function individually with cProfile (Section 1), then measure end-to-end TTFA in the full parallel pipeline (Section 2). Finally, we benchmark the Channel throughput against \texttt{threading.Queue} (Section 3).

\paragraph{Test Setup}
\begin{itemize}
    \item \textbf{Audio input}: Synthetic speech generated via macOS \texttt{say} (4.4s, 48kHz mono)
    \item \textbf{ASR}: Groq Whisper API (\texttt{whisper-large-v3-turbo})
    \item \textbf{LLM}: Groq \texttt{llama-3.3-70b-versatile} (default), also tested OpenAI \texttt{gpt-5.4}
    \item \textbf{TTS}: Inworld TTS API (\texttt{inworld-tts-1.5-max})
\end{itemize}

\subsection{cProfile Analysis (per-function)}

Each function was profiled individually with its own \texttt{cProfile} session, running components one at a time and feeding output $\rightarrow$ input sequentially.

\subsubsection{\texttt{VAD.\_process\_audio\_frame} --- 219 calls, 81ms total (0.4ms/call)}

The VAD processes every 20ms audio chunk. The hot path is Silero VAD inference:

\begin{table}[h]
\centering
\begin{tabular}{|l|r|l|}
\hline
\textbf{Function} & \textbf{Cumulative Time} & \textbf{Notes} \\
\hline
\texttt{torch.nn.Module.\_call\_impl} & 60 ms & Silero VAD neural network forward pass \\
\texttt{AudioFrame.get} (resampling) & 4 ms & numpy \texttt{interp} for 48kHz$\rightarrow$16kHz \\
\texttt{torch.tensor} & 4 ms & Converting audio chunks for Silero \\
\texttt{deque.popleft} & 3 ms & VAD buffer management \\
\hline
\end{tabular}
\end{table}

In the pipeline, Smart Turn ONNX inference adds $\sim$13ms per chunk during silence detection, but the core VAD loop itself is very fast.

\subsubsection{\texttt{ASR.\_transcribe\_audio} --- 1 call, 553ms}

Dominated entirely by the Groq Whisper HTTP API call:

\begin{table}[h]
\centering
\begin{tabular}{|l|r|l|}
\hline
\textbf{Function} & \textbf{Cumulative Time} & \textbf{Notes} \\
\hline
\texttt{requests.post} $\rightarrow$ \texttt{sessions.send} & 550 ms & Full HTTP round-trip \\
\texttt{ssl.\_SSLSocket.read} & 465 ms & Waiting for Groq server response \\
\texttt{urllib3.connection.connect} & 72 ms & TLS handshake \\
\texttt{\_prepare\_audio\_for\_transcription} & $\sim$3 ms & WAV file creation + resampling \\
\hline
\end{tabular}
\end{table}

99\% of time is network I/O.

\subsubsection{\texttt{LLM} TTFT (\texttt{litellm.completion} $\rightarrow$ first token) --- 540ms}

cProfile captures the full \texttt{litellm.completion()} call up to the first streaming token:

\begin{table}[h]
\centering
\begin{tabular}{|l|r|l|}
\hline
\textbf{Function} & \textbf{Cumulative Time} & \textbf{Notes} \\
\hline
\texttt{litellm.utils.wrapper} & 535 ms & litellm entry point \\
\texttt{token\_counter} $\rightarrow$ \texttt{from\_pretrained} & 312 ms & Tokenizer loading (one-time cost) \\
\texttt{litellm.main.completion} & 217 ms & Actual API call to Groq \\
\texttt{httpx.Client.send} & 198 ms & HTTP request + first response chunk \\
\hline
\end{tabular}
\end{table}

\textbf{Key finding}: 312ms (58\%) of the TTFT is litellm's tokenizer initialization (\texttt{from\_pretrained}), not the actual API call. On subsequent calls this is cached and the real TTFT drops to $\sim$200ms.

\subsubsection{\texttt{TTS} \texttt{requests.post} (Inworld API) --- 522ms}

Profiles the Inworld TTS API call including streaming audio response:

\begin{table}[h]
\centering
\begin{tabular}{|l|r|l|}
\hline
\textbf{Function} & \textbf{Cumulative Time} & \textbf{Notes} \\
\hline
\texttt{urllib3.connectionpool.urlopen} & 492 ms & Full HTTP round-trip + streaming \\
\texttt{ssl.\_SSLSocket.read} & 370 ms & Reading streaming audio chunks \\
\texttt{http.client.\_read\_status} & 368 ms & Waiting for first response byte \\
\texttt{urllib3.connection.connect} & 122 ms & TLS handshake to Inworld \\
\texttt{create\_connection} & 74 ms & TCP socket setup \\
\hline
\end{tabular}
\end{table}

TLS handshake (122ms) is a significant fixed cost per request.

\subsubsection{\texttt{Channel.send} / \texttt{Channel.\_wait\_and\_get} --- 100K calls, 346ms total (3.5$\mu$s/call)}

cProfile of the send path in a 1P1C throughput test (100K messages):

\begin{table}[h]
\centering
\begin{tabular}{|l|r|l|}
\hline
\textbf{Function} & \textbf{Cumulative Time} & \textbf{Notes} \\
\hline
\texttt{Sender.send} & 262 ms & Entry point --- iterates channels, tracks metrics \\
\texttt{Channel.\_send} & 125 ms & Appends item + \texttt{notify\_all()} under Condition lock \\
\texttt{Channel.\_gc} & 56 ms & \texttt{min(cursors.values())} + \texttt{del items[:drop]} on every receive \\
\texttt{threading.notify\_all} & 40 ms & Wakes all waiting threads (Queue uses \texttt{notify} --- wake one) \\
\texttt{sys.getsizeof} & 8 ms & Byte count tracking on every send \\
\hline
\end{tabular}
\end{table}

The main overhead vs \texttt{threading.Queue}: \texttt{\_gc()} runs on every receive (dict scan + list slice), and \texttt{notify\_all()} wakes all threads instead of one. Queue's \texttt{put()}/\texttt{get()} is simpler: acquire lock $\rightarrow$ append to deque $\rightarrow$ notify one $\rightarrow$ release.

\subsection{End-to-End Pipeline TTFA}

After profiling each function individually, we ran the full pipeline in parallel to measure real TTFA:

\begin{center}
\texttt{FileSource $\rightarrow$ VAD $\rightarrow$ ASR $\rightarrow$ LLM $\rightarrow$ TTS $\rightarrow$ NullSink}
\end{center}

TTFA is measured from the moment VAD finalizes the user's speech segment (end-of-utterance) to the first TTS audio chunk arriving at the sink. Each run is a separate process to avoid caching effects.

\subsubsection{TTFA Breakdown (averaged over 10 pipeline runs)}

\begin{table}[h]
\centering
\begin{tabular}{|l|r|r|r|l|}
\hline
\textbf{Stage} & \textbf{Avg} & \textbf{Min} & \textbf{Max} & \textbf{What it measures} \\
\hline
ASR (speech-to-text) & 310 ms & 169 ms & 820 ms & VAD-end $\rightarrow$ ASR transcription complete \\
LLM TTFT (first token) & 613 ms & 453 ms & 732 ms & ASR-done $\rightarrow$ first LLM token streamed \\
TTS TTFB (first audio) & 406 ms & 305 ms & 589 ms & First LLM token $\rightarrow$ first TTS audio chunk \\
\textbf{TTFA} & \textbf{1,330 ms} & \textbf{1,033 ms} & \textbf{1,873 ms} & \textbf{VAD-end $\rightarrow$ first TTS audio (sum of above)} \\
\hline
\end{tabular}
\end{table}

\subsubsection{Pipeline Hop Analysis}

We confirmed there is no hidden overhead in the channel/thread infrastructure:

\begin{table}[h]
\centering
\begin{tabular}{|l|r|}
\hline
\textbf{Hop} & \textbf{Latency} \\
\hline
VAD finalize (segment concat) & 1 ms \\
VAD $\rightarrow$ ASR channel transit & 0 ms \\
ASR $\rightarrow$ Adapter $\rightarrow$ LLM channel transit & 0 ms \\
\hline
\end{tabular}
\end{table}

\subsubsection{LLM TTFT Comparison (Isolated)}

We measured raw LLM TTFT with no pipeline overhead:

\begin{table}[h]
\centering
\begin{tabular}{|l|r|r|r|}
\hline
\textbf{Model} & \textbf{Avg TTFT} & \textbf{Min} & \textbf{Max} \\
\hline
Groq \texttt{llama-3.3-70b-versatile} & 390 ms & 257 ms & 646 ms \\
OpenAI \texttt{gpt-5.4} & 875 ms & 458 ms & 1,566 ms \\
OpenAI \texttt{gpt-4.1-nano} & 1,157 ms & 924 ms & 1,437 ms \\
\hline
\end{tabular}
\end{table}

Groq is the fastest due to their custom LPU inference hardware.

\subsubsection{Improvement Suggestions}

\paragraph{1. \texttt{LLM} TTFT --- Warm up tokenizer (613ms, 46\% of TTFA)}

\textbf{Problem}: cProfile reveals 312ms of the 540ms is litellm loading the tokenizer via \texttt{from\_pretrained}.

\textbf{Improvement}: Call \texttt{litellm.completion()} once at startup so the tokenizer is cached before real requests.

\paragraph{2. \texttt{TTS} \texttt{requests.post} --- Pre-warm TLS connection (406ms, 31\% of TTFA)}

\textbf{Problem}: TLS handshake to Inworld costs 122ms per request. Each TTS sentence creates a new connection.

\textbf{Improvement}: Use \texttt{requests.Session()} with connection pooling to reuse the TLS session across requests.

\paragraph{3. \texttt{ASR.\_transcribe\_audio} --- Use streaming ASR (310ms, 23\% of TTFA)}

\textbf{Problem}: Each API call to Groq Whisper takes $\sim$550ms and processes the entire segment at once.

\textbf{Improvement}: Use a WebSocket-based ASR service (e.g., Deepgram) that transcribes incrementally as audio arrives.

\paragraph{4. \texttt{VAD.\_process\_audio\_frame} --- Reduce Smart Turn frequency ($<$5ms/chunk, minimal TTFA impact)}

\textbf{Problem}: \texttt{\_check\_smart\_turn()} runs ONNX inference ($\sim$13ms) on every audio chunk during silence detection.

\textbf{Improvement}: Only run Smart Turn every 200ms instead of every 20ms. 5x/second is sufficient for turn detection.

\subsubsection{Results After Implementing Improvements 1 \& 2}

We implemented the top two improvements:

\begin{enumerate}
    \item \textbf{LLM tokenizer warmup}: Added \texttt{setup()} to \texttt{LLM} that calls \texttt{litellm.completion()} once at startup with a dummy prompt.
    \item \textbf{TTS persistent session}: Replaced \texttt{requests.post()} with \texttt{self.\_session.post()} using a \texttt{requests.Session()}.
\end{enumerate}

\begin{table}[h]
\centering
\begin{tabular}{|l|r|r|l|}
\hline
\textbf{Stage} & \textbf{Before (avg)} & \textbf{After (avg)} & \textbf{Improvement} \\
\hline
ASR & 310 ms & 280 ms & -30 ms \\
LLM TTFT & 613 ms & 448 ms & \textbf{-165 ms (27\%)} \\
TTS TTFB & 406 ms & 427 ms & +21 ms (API variance) \\
\textbf{TTFA} & \textbf{1,330 ms} & \textbf{1,155 ms} & \textbf{-175 ms (13\%)} \\
\hline
\end{tabular}
\end{table}

\subsection{Channel Throughput Benchmark}

We benchmarked our custom \texttt{Channel}/\texttt{Sender}/\texttt{Receiver} (function 5) against Python's \texttt{threading.Queue} using methodology adapted from the LMAX Disruptor perftest suite: 100K messages per run, 7 runs per test, \texttt{gc.collect()} between runs, correctness checksums on every run.

\subsubsection{Throughput (median of 7 runs)}

\begin{table}[h]
\centering
\begin{tabular}{|l|r|r|l|}
\hline
\textbf{Topology} & \textbf{Channel} & \textbf{Queue} & \textbf{Ratio} \\
\hline
1P1C (unicast) & 512K ops/s & 1,244K ops/s & Queue 2.4x faster \\
1P3C (multicast) & 215K ops/s & 414K ops/s & Queue 1.9x faster \\
3P1C (fan-in) & 158K ops/s & 1,188K ops/s & Queue 7.5x faster \\
Pipeline (P$\rightarrow$S1$\rightarrow$S2$\rightarrow$S3) & 120K ops/s & 405K ops/s & Queue 3.4x faster \\
\hline
\end{tabular}
\end{table}

\subsubsection{Latency (ping-pong round-trip, best of 5 runs)}

\begin{table}[h]
\centering
\begin{tabular}{|l|r|r|}
\hline
\textbf{Percentile} & \textbf{Channel} & \textbf{Queue} \\
\hline
p50 & 8.4 $\mu$s & 9.0 $\mu$s \\
p90 & 10.5 $\mu$s & 9.6 $\mu$s \\
p99 & 13.8 $\mu$s & 12.0 $\mu$s \\
p99.9 & 22.5 $\mu$s & 23.3 $\mu$s \\
max & 29.7 $\mu$s & 103.0 $\mu$s \\
\hline
\end{tabular}
\end{table}

Channel has slightly lower median latency and better tail (30$\mu$s max vs 103$\mu$s).

\subsubsection{Improvement: Rust Ring Buffer with CAS (FastChannel)}

\textbf{Problem}: Channel is 2--7x slower than Queue on throughput. cProfile identifies \texttt{\_gc()} on every receive, \texttt{notify\_all()} on every send, and \texttt{sys.getsizeof()}/\texttt{time.time()} per-message metrics as the main overhead.

\textbf{Improvement}: Replace Channel's Python \texttt{list} + \texttt{threading.Condition} with a fixed-size ring buffer implemented in Rust via PyO3. The producer advances a write sequence counter via atomic \texttt{fetch\_add} (CAS), and consumers advance their cursors via \texttt{compare\_exchange} (CAS) --- no mutex on the hot path. The ring buffer eliminates per-message GC (old slots are simply overwritten when the ring wraps).

\subsubsection{Results After Implementing FastChannel}

\begin{table}[h]
\centering
\begin{tabular}{|l|r|r|r|}
\hline
\textbf{Topology} & \textbf{Channel} & \textbf{Queue} & \textbf{FastChannel} \\
\hline
1P1C (unicast) & 512K ops/s & 1,244K ops/s & \textbf{2,921K ops/s} \\
1P3C (multicast) & 215K ops/s & 414K ops/s & \textbf{1,798K ops/s} \\
3P1C (fan-in) & 158K ops/s & 1,188K ops/s & \textbf{2,664K ops/s} \\
Pipeline (P$\rightarrow$S1$\rightarrow$S2$\rightarrow$S3) & 120K ops/s & 405K ops/s & \textbf{981K ops/s} \\
\hline
\end{tabular}
\end{table}

FastChannel is \textbf{2--4x faster than Queue} and \textbf{5--8x faster than the original Channel}.

\begin{table}[h]
\centering
\begin{tabular}{|l|r|r|r|}
\hline
\textbf{Percentile} & \textbf{Channel} & \textbf{Queue} & \textbf{FastChannel} \\
\hline
p50 & 8.4 $\mu$s & 9.0 $\mu$s & \textbf{7.5 $\mu$s} \\
p90 & 10.5 $\mu$s & 9.6 $\mu$s & \textbf{8.0 $\mu$s} \\
p99 & 13.8 $\mu$s & 12.0 $\mu$s & \textbf{10.4 $\mu$s} \\
p99.9 & 22.5 $\mu$s & 23.3 $\mu$s & \textbf{17.6 $\mu$s} \\
max & 29.7 $\mu$s & 103.0 $\mu$s & \textbf{45.6 $\mu$s} \\
\hline
\end{tabular}
\end{table}

FastChannel has the lowest latency at every percentile.
